\documentclass[12pt]{article}
\usepackage[utf8]{inputenc}
\usepackage[T1]{fontenc}
\usepackage{amsmath, amssymb, geometry}
\usepackage{enumitem}
\geometry{margin=0.85in}
\setlength{\parindent}{0pt}
\renewcommand{\arraystretch}{1.2}

\title{Quiz 5\\ \vspace{5pt} \Large ECON 308 -- International Economic Relations}
\author{Instructor: Muhebullah Karimzada}
\date{December 02,  2025}
\usepackage{comment}

\begin{document}
\maketitle

\section*{Multiple-Choice Questions (1 mark each)}



\begin{enumerate}[label=\textbf{\arabic*.}]

\item A \textit{specific tariff} is defined as:
\begin{enumerate}[label=\Alph*.]
\item A tax levied as a percentage of the value of the imported good
\item A tax levied as a fixed amount per unit of the imported good
\item A tax that varies depending on domestic production
\item A tax imposed only on exports
\end{enumerate}
\vspace{0.1 cm}

\item For a \textit{small country}, the welfare effect of a tariff includes:
\begin{enumerate}[label=\Alph*.]
\item Terms-of-trade gain only
\item Terms-of-trade gain minus deadweight loss
\item Only deadweight loss
\item Deadweight loss plus quota rents
\end{enumerate}

\vspace{0.1 cm}
\item A quota and a tariff are equivalent when:
\begin{enumerate}[label=\Alph*.]
\item Demand is perfectly inelastic
\item The foreign government collects the quota rents
\item The quota licenses are auctioned competitively
\item The tariff revenue is negative
\end{enumerate}
\vspace{0.1 cm}

\item A \textit{voluntary export restraint (VER)} differs from a quota because:
\begin{enumerate}[label=\Alph*.]
\item It applies only to exports from the importing country
\item The importing government collects the rents
\item The exporting country collects the rents
\item It has no welfare effects
\end{enumerate}

\item A \textit{large country} imposing a tariff:
\begin{enumerate}[label=\Alph*.]
\item Cannot affect world prices
\item May improve its terms of trade
\item Always reduces its welfare
\item Always improves foreign welfare
\end{enumerate}
\vspace{0.1 cm}
\item The two types of deadweight loss from a tariff are:
\begin{enumerate}[label=\Alph*.]
\item Government revenue and quota rents
\item Production distortion and consumption distortion
\item Producer surplus loss and foreign revenue loss
\item Terms-of-trade loss and foreign welfare loss
\end{enumerate}
\vspace{0.1 cm}

\item The \textit{effective rate of protection} (ERP) measures:
\begin{enumerate}[label=\Alph*.]
\item The tariff rate on the final good
\item The average tariff rate across all industries
\item The percentage increase in value added due to tariffs
\item The tariff burden on foreign exporters
\end{enumerate}

\end{enumerate}

\vspace{1em}

\vspace{1em}



\subsection*{Problem 1 (13 marks)}

The domestic market for a good is described by:
\[
Q_S = P - 50, \qquad Q_D = 400 - P.
\]
The world price is \( P_W = 150 \).  
The country is \textit{small}.

\begin{enumerate}[label=(\arabic*)]
\item Determine the free-trade domestic price, quantity supplied, quantity demanded, and imports.
\item The government imposes a \textbf{specific tariff} of \( t = 30 \).  
Compute the new domestic price, \( Q_S \), \( Q_D \), and imports.
\item Calculate government \textbf{tariff revenue}.
\item Compute the \textbf{deadweight loss} of the tariff (production + consumption distortions).
\end{enumerate}



\end{document}
